Käesoleva töö eesmärk on luua alus eestikeelse äratussõna tuvastuse realiseerimiseks piiratud ressursiga nutikodu mikrokontrolleril. Seni tehtud töö põhjal on selgunud, et lisaks mudeli treenimisele tuleb sama tõsiselt võtta ka hindamismetoodikat ja andmestiku rollide eristamist.

Töö käigus rekonstrueeriti varasemad ebaõnnestunud katsed, parandati \texttt{microWakeWord} põhist treeningu- ja raportitoru ning valideeriti see avaliku \texttt{Speech Commands} andmestiku peal. Selle tulemusena on treeningutoru reprodutseeritavalt kontrollitud \texttt{Speech Commands} andmestikul ning valmis eestikeelsetele andmetele ümber suunamiseks.

Peamine vahekokkuvõte on, et äratussõna lahenduse kvaliteeti tuleb hinnata mitte ainult positiivsete ja negatiivsete klippide, vaid ka pika ambient-signaali peal. See eristus on töö järgmiste etappide hindamismetoodika alus, mis peab aitama hinnata valevallandumiste riski reaalses kasutuses.
